\section{Methods}\label{sec6}

\subsection{Study design and analytic scope}

This study analyses naturalistic public human--LLM conversations as observational digital traces of everyday work and study. We use three public conversational corpora: WildChat-4.8M (public non-toxic release)~\cite{allenai2025wildchat48m} (reported below as WildChat), LMSYS Chat-1M~\cite{zheng2023lmsyschat1m} (reported below as LMSYS Chat) and the ShareChat~\cite{yan2025sharechat} strict-English subset (reported below as ShareChat). Within these corpora, we focus on two common task ecologies, coding and writing, because both are frequent everyday LLM uses but expose different opportunities for visible learning-oriented behaviour. Coding conversations often contain debugging, error diagnosis and iterative testing; writing conversations more often contain drafting, revision and directive editing~\cite{keuning2018systematic,flower1981cognitive,sommers1980revision}.

The design is descriptive and associational, following the logic of large-scale digital trace and computational social-science analyses~\cite{lazer2009computational,salganik2017bitbybit}. The empirical sequence mirrors the Results section. We first describe whether learning-oriented engagement is visible in ordinary use, then ask where constructive engagement is concentrated, whether scaffolded assistant support marks richer constructive participation, which support forms carry that association and how support and engagement are ordered at the adjacent-turn scale. These analyses support claims about observable behavioural signatures in conversation logs. They do not estimate randomized exposure effects, durable learning outcomes, changes in users' knowledge or stable learner traits.

\subsection{Corpus construction and task settings}

The analytic corpus contains 128,569 conversations and 981,470 total turns, including 491,685 user turns and 489,785 assistant turns. The six task settings are WildChat coding, WildChat writing, LMSYS Chat coding, LMSYS Chat writing, ShareChat coding and ShareChat writing. WildChat and LMSYS Chat provide large public human--LLM interaction logs from deployed chat systems, whereas ShareChat contributes a smaller strict-English analytic subset after our language screen. The settings are complementary naturalistic environments rather than experimentally matched platform conditions~\cite{zhao2024wildchat,zheng2023lmsyschat1m,yan2025sharechat}.

Corpus construction followed the same high-level sequence across datasets. Public source records were converted to role-ordered user--assistant conversations; conversations were restricted to English-oriented exchanges with at least four turns; coding and writing task ecologies were assigned with LLM-assisted semantic filters; and retained conversations were passed through the shared annotation pipeline. When keyword prefilters were used, the LLM-assisted semantic task filter retained conversations with a score of at least 0.70. ShareChat additionally underwent a strict-English screen after task filtering, retaining only conversations whose final conversation-level language flag was English. Supplementary Tables~\ref{tab:corpus_provenance_artifacts} and \ref{tab:filter_classifier_validation} provide the raw-to-final provenance ledger, classifier sources and validation evidence for these filtering decisions.

\subsection{Operational constructs}

\subsubsection{Task ecology, user framing and depth}

Each conversation is represented as an ordered sequence of user and assistant turns together with available metadata. The main contextual variables are data source, task ecology, user framing, topic and language. Task ecology and user framing are analytic variables generated by the filtering and annotation pipeline; topic, language and model/source metadata are retained where available. User framing is a conversation-level LLM-assisted label for whether the opening user request or early exchange explicitly signals understanding, exploration or learning, as opposed to a request without such explicit learning-oriented framing. We interpret this as an observed textual signal, not as a stable measure of motivation or learner identity.

Interaction depth is measured from the number and ordering of turns in the conversation. It is used descriptively, as a factor of interactional opportunity and sustained exchange, not as an independent manipulation. Post-answer depth is measured as the number of turns after the first assistant response. Because any-event outcomes are easier to observe in longer conversations, the supplementary analyses include rate-based sensitivity checks for constructive engagement.

\subsubsection{User engagement}

The unit of annotation is the conversational turn. User turns are labelled for behavioural, emotional and cognitive engagement, drawing on engagement theory and the ICAP distinction between passive, active and constructive participation~\cite{fredricks2004school,chi2014icap}. Passive engagement (\texttt{P}) captures reception or acknowledgement without substantial transformation. Active engagement (\texttt{A}) captures task-relevant action, follow-up or manipulation of supplied information. Constructive engagement (\texttt{C}) captures deeper work with ideas, such as testing an implication, revising a hypothesis, articulating a constraint, transferring a mechanism or extending an explanation. Visible emotional expression is labelled separately and reported descriptively because it was rare.

The primary user-side outcome is constructive engagement, treated as the highest-level behavioural proxy for learning-oriented participation available in these logs. At the turn level, constructive engagement is a binary indicator for whether the user turn is labelled \texttt{C}. At the conversation level, we use both the count of constructive user turns and an indicator for whether a conversation contains at least one constructive user turn. Constructive ratios are defined as constructive user turns divided by user turns.

\subsubsection{Assistant support}

Assistant turns are labelled for scaffolding. Scaffolded support denotes assistant behaviour that structures, regulates, diagnoses or redirects the user's work while leaving some responsibility with the user; the comparison category is a non-scaffolded reference response, including straightforward task fulfilment or stand-alone information provision. This definition follows the scaffolding tradition, especially contingency, transfer of responsibility and potential fading~\cite{wood1976role,van2010scaffolding,quintana2018scaffolding}.

The primary assistant-side exposure is scaffolding, operationalized as scaffolded support. At the turn level, this is the distinction between scaffolded support and a non-scaffolded reference response. At the conversation level, we use an indicator for whether a conversation contains at least one scaffolded assistant turn. Scaffolded turns can also receive support-intent labels (\texttt{I1} metacognitive, \texttt{I2} cognitive and \texttt{I3} affective support) and non-mutually exclusive support-form labels (\texttt{M1} feedback, \texttt{M2} hinting, \texttt{M3} instructing, \texttt{M4} explaining, \texttt{M5} modelling and \texttt{M6} questioning). Support-form analyses condition on scaffolded turns and compare conversations containing each support-form label with scaffolded conversations that do not contain that label. Operational definitions are provided in Supplementary Table~\ref{tab:annotation_codebook_definitions}; paraphrased examples and two complete de-identified case transcripts are provided in Supplementary Tables~\ref{tab:label_examples}--\ref{tab:complete_writing_case}.

\subsection{Annotation workflow and validation}

The annotation workflow combined theory-grounded codebook development, prompt iteration, parser checks, targeted post-processing, human review and production-scale LLM-assisted labelling. We used LLM-assisted annotation as a scalable text-analysis workflow while retaining human verification because prior work shows both the promise and the need for validation when language models are used as annotators~\cite{gilardi2023chatgpt,ziems2024can}. After the workflow stabilized, the same construct definitions, schema checks and post-processing logic were applied to the WildChat, LMSYS Chat and ShareChat analytic corpora. The archived annotation metadata record the provider, deployment family, API version, prompt version, retry policy, parser checks and rule-based post-processing rules used for production labelling; Supplementary Section A.3 gives the reproducibility details.

Human verification is reported in layers. Human--human agreement assesses whether trained reviewers could apply the codebook consistently. Human--LLM production-label checks compare final or post-processed production labels with human-confirmed labels. Targeted post-production audits check whether narrow rule-based reviews reduced known residual false positives. Agreement is summarized with per-label F1, Matthews correlation coefficient and Gwet's AC1 because the labels are sparse, prevalence-imbalanced and sometimes non-mutually exclusive~\cite{matthews1975comparison,gwet2008computing}. The consolidated production-label audits indicate acceptable or better reliability across the labels used for the main claims: intentional framing reached F1=0.857, constructive-versus-non-constructive user engagement reached F1=0.922, broad scaffolded support reached F1=0.912 and support-form labels ranged from F1=0.691 to 0.938. Supplementary Tables~\ref{tab:label_validation} and \ref{tab:production_label_validation} report the complete validation metrics and audit designs.

\subsection{Statistical analyses}

Descriptive analyses estimate engagement prevalence, scaffolding prevalence, passive/active/constructive composition, visible emotional-expression rates and behavioural-depth summaries within each task setting. Context analyses compare observed user framing, task ecology and depth groups. Because the probability that a conversation contains at least one constructive turn partly reflects the number of turns available for observing one, we also fit a grouped-binomial constructive-rate sensitivity model in which the outcome is the constructive user-turn count out of total user turns, with user framing, task ecology, length bucket and dataset fixed effects as predictors.

Conversation-level support associations compare conversations with and without at least one scaffolded assistant turn. Constructive-ratio contrasts use turn-weighted ratios within each comparison group, with uncertainty from 1,000 conversation-level bootstrap resamples. Adjusted Section~2.3 models use setting-specific Poisson generalized linear models for raw constructive-turn counts and companion logistic models for whether a conversation contains at least one constructive user turn. The primary Poisson model treats conversation length as a covariate rather than using a user-turn offset; offset-rate and quasi-Poisson sensitivities are reported in the Supplement. These models are interpreted as adjusted associations, not causal effects.

Support-form contrasts compare scaffolded conversations containing a given support form with scaffolded conversations not containing that form. Because support-form labels are non-mutually exclusive and nested within scaffolded support, these contrasts are descriptive form-pattern analyses rather than mutually exclusive treatment comparisons. Bracket tests in Fig.~\ref{fig:support_form_supply} use Benjamini--Hochberg false-discovery-rate adjustment within each displayed family of six support-form comparisons.

Temporal analyses focus on adjacent-turn ordering. We compare the probability that the next user turn is constructive after scaffolded versus non-scaffolded assistant turns, estimate the same contrast within prior user states and estimate the probability that the next assistant turn is scaffolded after different prior user states. Adjacent-turn contrasts use 1,000 conversation-cluster bootstrap resamples because multiple turn pairs can come from the same conversation. Integrated adjacent-turn regressions combine user-state features, assistant-scaffolding features and adjacent-turn context in one model, with conversation-clustered standard errors. Supplementary Tables~\ref{tab:model_specs}--\ref{tab:temporal_estimands} and \ref{tab:setting_model_covariates} give the complete estimand map, covariate sets, model specifications and sensitivity checks.

\subsection{LLM use and reporting boundaries}

LLMs were used as objects of study and as annotation tools. The authors also used LLM-based tools to assist with code drafting, debugging, analysis workflow organization and language editing. LLM-generated assistance was not used as a substitute for human verification of the validation samples or for interpretation of the statistical results.

All reported results are aggregate turn-level or conversation-level analyses. The study does not link conversations into user histories, estimate user-level longitudinal trajectories or redistribute verbatim raw conversation logs. Operational examples are de-identified and paraphrased or lightly edited, and temporal analyses are confined to ordering within a conversation. Model/source checks are used as robustness analyses rather than model-causal comparisons, because model labels in public logs are not experimentally assigned and may be confounded with time, user population, task mix and logging context.
